\documentclass{article}
\usepackage[utf8]{inputenc}
\usepackage[T1]{fontenc}
\usepackage[english]{babel}
\usepackage{amsmath,amssymb,amsthm}
\usepackage{graphicx}
\usepackage{listings}
\usepackage{tikz}
\usetikzlibrary{matrix,arrows.meta,positioning,calc,fit}
\usepackage{tikz-cd}
\newsavebox{\diagbox}
\usepackage{hyperref}

\theoremstyle{definition}
\newtheorem{definition}{Definition}[section]
\newtheorem{example}[definition]{Example}
\newtheorem{remark}[definition]{Remark}

\theoremstyle{plain}
\newtheorem{theorem}[definition]{Theorem}

\lstset{
  basicstyle=\tiny\ttfamily,
  breaklines=true,
  columns=fullflexible,
  frame=single,
  showstringspaces=false,
  keepspaces=true
}

\title{Architecture of Mathematics:\\
Internal Languages, Categorical Models, \\
and the Spectral Lattice of Formal Systems}
\author{Groupoid Infinity}
\date{\today}

\begin{document}

\maketitle

\begin{abstract}
In this article, we investigate the conjecture that no single formal language can serve effectively as a universal
foundation for the entirety of mathematics, proposing instead that mathematics is best served by a
system of specialized languages connected by functorial transports. We analyze the Groupoid Infinity
stack of formal languages as a \emph{lattice} of internal languages, rather than a linear tower.
The languages fall into three distinct branches: a \emph{constructive type theory and set analysis}
spine (Henk--Frank--Christine), a \emph{super homotopy} spine extending to differential cohesion
and quantum spectra (Henk--Per--Anders--Urs), and an \emph{infinity category theory} spine (Dan--Mike--Ulrik).
Due to their divergent semantic and syntactic structures, attempting to unify them into a single monolithic
abstract syntax tree (AST) is often impractical and conceptually strained. We demonstrate how these separate
domains interact functorially --- for instance, simplicial type theory (Ulrik/STT) provides the
directed-homotopy backbone for Urs's 3-layer architecture. Finally, we propose a cohomological perspective
on language semantics: by associating syntactic CW-complexes and chain complexes to individual languages,
one might define their cohomology groups and assemble a spectral sequence over the full lattice
of language categories, outlining a conceptual mathematical framework that formalizes this ecosystem as a
Grothendieck fibration of syntactic categories under Lawvere's functorial semantics.
\end{abstract}

\newpage
\tableofcontents
\newpage

\section{Categorical Foundations}

\begin{definition}[Categories with Families (CwF)]
A CwF consists of a category $\mathcal{C}$ of contexts and substitutions with a terminal object, equipped with a functor $Ty : \mathcal{C}^{\mathrm{op}} \to \mathbf{Set}$, a presheaf of terms $Tm : (\int Ty)^{\mathrm{op}} \to \mathbf{Set}$, and a context comprehension operation with projection morphisms.
\end{definition}

\begin{definition}[Natural Models]
Following Awodey, a natural model is a category $\mathcal{C}$ with a terminal
object $\mathbf{1}$ and a representable natural transformation $p : Tm \to Ty$ between presheaves.
Pullbacks of $p$ along Yoneda provide context extension and substitution.
\end{definition}

\begin{definition}[Grothendieck Fibration]
A functor $p : \mathcal{E} \to \mathcal{B}$ is a fibration when for every
$X \in \mathcal{E}$ and $u : J \to p(X)$ in $\mathcal{B}$ there exists a
Cartesian lift $f : Y \to X$ in $\mathcal{E}$.
\end{definition}

The category of languages $\mathcal{L}$ has formal languages as objects and
syntactic functors (compilers, type checkers, normalizers, extractors) as
morphisms.

\begin{table}[ht!]
\centering
\caption{Semantic Categories and Core Language Primitives}
\label{tab:categories_internal_languages}
\noindent\makebox[\textwidth][c]{%
\begin{tabular}{lll}
\hline
\textbf{Language} & \textbf{Semantic Category / Model} & \textbf{Application \& Physical/Mathematical Domain} \\
\hline
$O_\lambda$ & Cartesian Closed Category ($\mathbf{CCC}$) & Untyped $\lambda$-calculus (evaluation and interpretation) \\
$O_\pi$ & Symmetric Monoidal Category & Process calculi (CCS, CSP, Milner's $\pi$-calculus) \\
$O_\mu$ & Tensor Category & Tensor calculus (vectorization) \\
\hline
Henk ($O_\Pi$) & $\mathbf{\Pi}\text{-}\mathbf{Cat}$ (C-systems / CwF) & PTS Calculus of Constructions (functional completeness) \\
Frank & $\mathbf{\Pi}\text{-}\mathbf{Cat} + \mathbf{Ind}$ & Minimal Calculus of Inductive Constructions (strictly positive schemes) \\
Christine ($O_\Sigma$) & $\mathbf{LCCC} + \mathbf{Ind}$ & Full CIC (MLTT-72 contextual completeness, dependent sums) \\
Laurent & $\mathbf{Top} / \mathbf{Met}$ & Calculus and Functional Analysis (topological and metric spaces) \\
Per ($O_W$) & $\mathbf{\Pi W}\text{-}\mathbf{Pretopos}$ & Martin-L\"{o}f Type Theory MLTT-80 (inductive W-types) \\
Anders ($O_I$) & $\infty\text{-}\mathbf{Cat}$ & Cubical Type Theory with Path types, no HITs ($\infty$-categories) \\
Anders\text{+} & $\infty\text{-}\mathbf{Topos}$ & Cubical Type Theory with HITs (lightweight core, univalence) \\
\hline
Felix ($O_{\dashv}$) & $\mathbf{Cohesive}\ \infty\text{-}\mathbf{Topos}$ & Modal HoTT (synthetic differential supergeometry, flat/sharp) \\
Jack & $\mathbf{Diff}\ \infty\text{-}\mathbf{Topos}$ & Stable homotopy (spectra, differential forms, K-theory) \\
Urs ($O_H$) & $\mathbf{Urs}\ \infty\text{-}\mathbf{Topos}$ (Linear Types) & Quantum linear types (Hadamard operator, qubits, braids) \\
\hline
Dan ($O_{\triangleright}$) & $\mathbf{Pres}$ (Simplicial Sets / Chains) & Guarded recursion, operational algebra (face/degeneracy maps) \\
Mike & $\mathbf{1\text{-}Cat}$ & Directed 1-Category theory (directed paths, linear structures) \\
Ulrik & $\mathbf{(\infty, 1)\text{-}Cat}$ (Riehl-Shulman STT) & Simplicial $(\infty, 1)$-categories (Segal, directed interval) \\
\hline
\end{tabular}%
}
\end{table}

\newpage
\section{The Three Spines of the Language Lattice}
The collection of formal internal languages does not admit a totally ordered hierarchy. Rather, the architecture factors into three principal spines. The first spine (\textbf{Frank}--\textbf{Christine}--\textbf{Laurent}) is dedicated to constructive inductive schemes and functional analysis. The central spine (\textbf{Henk}--\textbf{Per}--\textbf{Anders}--\textbf{Urs}) models univalent homotopy types and cohesive $\infty$-topoi. The third spine (\textbf{Dan}--\textbf{Mike}--\textbf{Ulrik}) models directed synthetic category theory and discrete operational presentations. As formalized in the commutative diagram below, the respective syntaxes are mutually disjoint; all inter-language translations are strictly mediated by functorial morphisms, including left adjoint lifts, reflective embeddings, and forgetful coreflections.

\begin{center}
\begin{tikzcd}[row sep=large, column sep=large]
& \mathbf{Lambda\ Cube} \ar[d] & \\
& \mathbf{Henk} \ar[d, "\text{012$\Sigma$W}"] \ar[dl, bend right=15, "\text{Ind}"'] \ar[dr, bend left=15, "\text{Presentations}"] & \\
\mathbf{Frank} \ar[d, "\text{Id,Fix}"'] & \mathbf{Per} \ar[d, "\text{Path,Glue}"] \ar[l, "\text{W $\rightarrow$ Ind}"] \ar[r, bend left=15, "\text{Load}"] & \mathbf{Dan} \ar[d, dashed, "\text{Lift}"] \\
\mathbf{Christine} \ar[d, "\text{Analysis}"'] \ar[ur, dotted, <->, "\simeq"] \ar[r] & \mathbf{Anders} \ar[d, "\text{Modal}"] \ar[r, dashed, bend left=15, "\text{Forget}"] & \mathbf{Mike} \ar[d, dashed, "\text{Free}"] \\
\mathbf{Laurent} \ar[r, "\text{Smooth}"] \ar[dr, bend right=15, "\text{Forms}"'] & \mathbf{Jack} \ar[d, "\text{Diff}"] & \mathbf{Ulrik} \ar[ul, dashed, "\text{Core}"'] \ar[l, red, bend left=15, "\text{STT}"] \ar[dl, red, bend left=15, "\text{STT}"] \\
& \mathbf{Urs} \ar[d, "\text{Quantum}"] & \\
& \mathbf{Cosmic\ Cube} &
\end{tikzcd}
\end{center}

\begin{remark}[Christine $\simeq$ Per]
Christine (full CIC with Id, Ind, Fix, Prop) and Per (MLTT-80 with
$\Sigma$, Id, W, 0, 1, 2) are at a comparable level of expressiveness. Per
serves as the foundation for Anders but can be considered interchangeable
base systems (Ind for modeling general HIT as in cubicaltt). The choice between inductive schemes (Frank/Christine path) and
W-types (Per path) is a design decision motivated by purity, not a strict ordering.
\end{remark}

\newpage
\subsection{The Topoi Structure}
To analyze the gradual expansion of language primitives, we construct a fine-grained lattice of semantic categories radiating from the category of dependent products ($\mathbf{\Pi}\text{-}\mathbf{Cat}$ (Henk)). The system expands in three cardinal directions (adding dependent sums to yield locally cartesian closed categories $\mathbf{LCCC}$, inductive algebras $\mathbf{\Pi}\text{-}\mathbf{Cat} + \mathbf{Ind}$ (Frank), and bicartesian closed structures $\mathbf{BiCCC}$ (IPL)) before merging into rich categories of models such as $\mathbf{LCCC} + \mathbf{Ind}$ (Christine), $\mathbf{\Pi W}\text{-}\mathbf{Pretopos}$ (Per), $\infty\text{-}\mathbf{Grpd}$ (Anders), and cohesive $\infty$-topoi. On the right, the lattice incorporates the infinity category spine: operational algebra branch ($\mathbf{Pres}$ (Dan)), the directed 1-category branch ($\mathbf{1\text{-}Cat}$ (Mike)), and the simplicial $(\infty, 1)$-category branch ($\mathbf{(\infty, 1)\text{-}Cat}$ (Ulrik)). All cross-branch connections are detailed in the global map in Section 2.

\begin{theorem}[Universal Primitive Lattice]
The formal internal languages form a connected lattice structure across three spines, linked by the modal, synthetic, and directed structural morphisms shown below:
\end{theorem}

\begin{lrbox}{\diagbox}
\scriptsize
\begin{tikzcd}[row sep=large, column sep=0.0em]
& \mathbf{\Pi W}\text{-}\mathbf{Pretopos}\ \text{(Per)} & \mathbf{LCCC} \ar[l, dashed] & \mathbf{Pres}\ \text{(Dan)} \ar[d, dashed, "\text{Lift}"] \\
\infty\text{-}\mathbf{Grpd}\ \text{(Anders)} & \mathbf{\Pi}\text{-}\mathbf{Cat}\ \text{(Henk)} \ar[ur] \ar[dr] \ar[l, dashed] \ar[r] & \mathbf{BiCCC}\ \text{(IPL)} \ar[r, dashed] & \mathbf{1\text{-}Cat}\ \text{(Mike)} \ar[d, dashed, "\text{Free}"] \\
\mathbf{Cohesive}\ \infty\text{-}\mathbf{Topos}\ \text{(Felix)} \ar[u, leftarrow, "\text{Modal}"'] & \mathbf{LCCC} + \mathbf{Ind}\ \text{(Christine)} & \mathbf{\Pi}\text{-}\mathbf{Cat} + \mathbf{Ind}\ \text{(Frank)} \ar[l, dashed] & \mathbf{(\infty, 1)\text{-}Cat}\ \text{(Ulrik)} \\
\mathbf{Diff}\ \infty\text{-}\mathbf{Topos}\ \text{(Jack)} \ar[u, leftarrow, "\text{Diff}"'] & & & \\
\mathbf{Urs}\ \infty\text{-}\mathbf{Topos}\ \text{(Urs)} \ar[u, leftarrow, "\text{Urs}"'] & & &
\end{tikzcd}%
\end{lrbox}
\noindent\makebox[\textwidth][c]{\usebox{\diagbox}}

\newpage
\subsection{The Semantic Foundations: Structural, Cohesive, and Directed Modalities}
The lattice stratifies into structural, cohesive, linear, and directed semantics. The foundational correspondences between semantic models and their internal syntaxes progressively build up structure:

\begin{itemize}
    \item \textbf{Structural Core:} At the base, Locally Cartesian Closed Categories ($\mathbf{LCCC}$) provide the semantics for $\mathbf{MLTT}$, the pure structural type theory \cite{Lof84}. Upgrading the semantics to $\infty$-categories adds globular path equality, yielding Cubical Type Theory with Path types but without Higher Inductive Types (HITs) \cite{cchm}. Finally, a full $\infty$-topos models Cubical Type Theory with HITs, which maintains a very lightweight core and syntax.
    \item \textbf{Homotopical Presheaves:} Bridging the gap from strict 1-categories to full $\infty$-categories, Thierry Coquand's Groupoid-Valued Presheaf models ($\mathbf{Grpd}\text{-}\mathbf{PSh}$) \cite{Coq88} provide a concrete constructive semantics for univalence, where dependent types are interpreted via the \emph{Grothendieck Construction}. This establishes the semantic foundation for Cubical Type Theory.
    \item \textbf{Cohesive Modalities:} Cohesive Type Theory (Felix) \cite{schreibershulman12cohesive} is developed specifically to model modalities such as connectedness, compactness, and infinitesimal shapes. Differential Cohesive Type Theory acts as the internal language for Differential Cohesive $\infty$-topoi \cite{schreiber13differential}, adding extra structure to $\infty$-topoi just as $\infty$-categories add path equality to LCCCs. These are typically modeled in the base library via undefined axioms (see \texttt{cohesivett}).
    \item \textbf{Linear/Tensor Modalities:} Symmetric Monoidal Categories ($\mathbf{SMC}$) supply the semantics for resource-sensitive and directed computations. Here, Process Calculus and Interaction Nets (Lafont) act as the internal language, representing modality types with special comonadic \emph{spawn} arrows that hide the underlying run-time implementation.
    \item \textbf{Directed and Simplicial Topoi:} Moving toward spatial and temporal operational algebras, presheaf categories ($\mathbf{Pres}$) provide the semantics for Guarded Type Theory (Dan) \cite{birkedal2011first}. Incorporating strict directional structures leads to $\mathbf{1\text{-}Cat}$, acting as the semantics for Directed Type Theory (Mike) \cite{Shulman15}. Finally, scaling to higher dimensions, $\mathbf{(\infty, 1)\text{-}Cat}$ serves as the model for Simplicial Type Theory (Ulrik / Riehl-Shulman STT) \cite{riehl2017type, gratzer2025yoneda}, extending directed paths with a Segal-like directed interval.
\end{itemize}

\newpage
This establishes a rich hierarchy of correspondences:

\begin{center}
\begin{tikzcd}[row sep=large, column sep=huge]
\mathbf{LCCC} \ar[r, leftrightarrow, "\text{Internal Language}"] & \mathbf{MLTT} \\
\mathbf{\Pi W}\text{-}\mathbf{Pretopos} \ar[r, leftrightarrow, "\text{Internal Language}"] & \mathbf{MLTT\ (W-types)} \\
\infty\text{-}\mathbf{Cat} \ar[r, leftrightarrow, "\text{Internal Language}"] & \mathbf{Cubical\ (Paths,\ no\ HITs)} \\
\infty\text{-}\mathbf{Topos} \ar[r, leftrightarrow, "\text{Internal Language}"] & \mathbf{Cubical\ (with\ HITs)} \\
\mathbf{Cohesive}\ \infty\text{-}\mathbf{Topos} \ar[r, leftrightarrow, "\text{Internal Language}"] & \mathbf{Cohesive\ HoTT} \\
\mathbf{SMC} \ar[r, leftrightarrow, "\text{Internal Language}"] & \mathbf{Process\ Calculus\ /\ Interaction\ Nets} \\
\mathbf{Pres} \ar[r, leftrightarrow, "\text{Internal Language}"] & \mathbf{Guarded\ Type\ Theory\ (Dan)} \\
\mathbf{1\text{-}Cat} \ar[r, leftrightarrow, "\text{Internal Language}"] & \mathbf{Directed\ Type\ Theory\ (Mike)} \\
\mathbf{(\infty,1)\text{-}Cat} \ar[r, leftrightarrow, "\text{Internal Language}"] & \mathbf{Simplicial\ Type\ Theory\ (Ulrik)}
\end{tikzcd}
\end{center}

\newpage
\section{Spine I: Set Analysis}
This spine represents the constructive/inductive spine of type theories and functional analysis. It is characterized by inductive algebraic constructions and real/functional analysis foundations.

\subsection{Frank: Minimal CIC \cite{Pfenning89}}\footnote{\url{https://frank.groupoid.space}}
Extends Henk with strictly positive inductive schemes.
\begin{table}[ht!]
\centering
\begin{tabular}{rl}
 cosmos :=&$\textbf{U}_j$ \\
    var :=&$\textbf{var}\ name$ \\
     pi :=&$\Pi\ name\ E\ E\ |\ \lambda\ name\ E\ E\ |\ E\ E$ \\
    ind :=&$\textbf{inductive}\ name\ \textrm{params}\ \textrm{constrs}\ |\ \textbf{constr}\ i\ \textrm{name}\ \textrm{args}\ |\ \textbf{ind}\ \textrm{name}\ E\ \textrm{args}\ E$ \\
\end{tabular}
\end{table}

\subsection{Christine: Full CIC (CwF/CwA) \cite{paulinmohring}}\footnote{\url{https://christine.groupoid.space}}
Complete Calculus of Inductive Constructions with $\Sigma$, Id, and Ind.
\begin{table}[ht!]
\centering
\begin{tabular}{rl}
 cosmos :=&$\textbf{U}_j$ \\
    var :=&$\textbf{var}\ name$ \\
     pi :=&$\Pi\ name\ E\ E\ |\ \lambda\ name\ E\ E\ |\ E\ E$ \\
  sigma :=&$\Sigma\ name\ E\ E\ |\ (E,E)\ |\ E.1\ |\ E.2$ \\
     id :=&$\textbf{Id}\ E\ |\ \textbf{ref}\ E\ |\ \textbf{id$_J$}\ E$ \\
    ind :=&$\textbf{inductive}\ name\ \textrm{params}\ \textrm{constrs}\ |\ \textbf{constr}\ i\ \textrm{name}\ \textrm{args}\ |\ \textbf{ind}\ \textrm{name}\ E\ \textrm{args}\ E$ \\
\end{tabular}
\end{table}

\subsection{Laurent: Calculus and Functional Analysis \cite{schwartz50distributions}}\footnote{\url{https://laurent.groupoid.space}}
Laurent formalizes classical real and functional analysis, distribution theory, and measure theory. It provides first-class support for reals, complex numbers, measures, limits, and integration.
\begin{table}[ht!]
\centering
\begin{tabular}{rl}
   cosmos :=&$\mathbf{Prop}\ :\ \mathbf{U_0}\ :\ \mathbf{U_1}$ \\
      var :=&$\mathbf{var}\ \mathrm{ident}\ |\ \mathbf{hole}$ \\
   forall :=&$\forall\ \mathrm{ident}\ \mathrm{E}\ \mathrm{E}\ |\ \lambda\ \mathrm{ident}\ \mathrm{E}\ \mathrm{E}\ |\ \mathrm{E}\ \mathrm{E}$ \\
   exists :=&$\exists\ \mathrm{ident}\ \mathrm{E}\ \mathrm{E}\ |\ (\mathrm{E}, \mathrm{E})\ |\ \mathrm{E}\mathbf{.1}\ |\ \mathrm{E}\mathbf{.2}$ \\
     base :=&$\mathbb{N}\ |\ \mathbb{Z}\ |\ \mathbb{Q}\ |\ \mathbb{R}\ |\ \mathbb{C}\ |\ \mathbb{H}\ |\ \mathbb{O}\ |\ \mathbb{V}^n$ \\
      set :=&$\mathbf{Set}\ |\ \mathbf{SeqEq}\ |\ \mathbf{And}\ |\ \mathbf{Or}\ |\ \mathbf{Complement}\ |\ \mathbf{Intersect}\ |\ \mathbf{Power}\ |\ \mathbf{Closure}\ |\ \mathbf{Cardinal}$ \\
        q :=&$\mathbf{-/\hspace{-1mm}\sim}\ |\ \mathbf{Quot}\ |\ \mathbf{Lift_Q}\ |\ \mathbf{Ind_Q}$ \\
       mu :=&$\mathbf{mu}\ |\ \mathbf{Measure}\ |\ \mathbf{Lebesgue}\ |\ \mathbf{Bochner}$ \\
      lim :=&$\mathbf{Seq}\ |\ \mathbf{Sup}\ |\ \mathbf{Inf}\ |\ \mathbf{Limit}\ |\ \mathbf{Sum}\ |\ \mathbf{Union}$ \\
\end{tabular}
\end{table}

\newpage
\section{Spine II: Super Homotopy Type Theory}
This spine represents the central path of homotopical, cubical, and cohesive type systems, transitioning from pure dependent products to higher-dimensional geometry.

\subsection{Henk: Pure Type Systems (PTS) \cite{Henk93}}\footnote{\url{https://henk.groupoid.space}}
The initial object in LCC categories. Contains only $\Pi$ and universes.
\begin{table}[ht!]
\centering
\begin{tabular}{rl}
 cosmos :=&$\textbf{U}_j$ \\
    var :=&$\textbf{var}\ name$ \\
     pi :=&$\Pi\ name\ E\ E\ |\ \lambda\ name\ E\ E\ |\ E\ E$ \\
\end{tabular}
\end{table}

\subsection{Per: Martin-L\"of Type Theory (MLTT-80) \cite{Lof84}}\footnote{\url{https://per.groupoid.space}}
MLTT-80 with W-types, base types 0, 1, 2, and cubical primitives.
\begin{table}[ht!]
\centering
\begin{tabular}{rl}
 cosmos :=&$\textbf{U}_j \ |\ \textbf{V}_k$ \\
    var :=&$\textbf{var}\ name\ |\ \textbf{hole}$ \\
     pi :=&$\Pi\ name\ E\ E\ |\ \lambda\ name\ E\ E\ |\ E\ E$ \\
  sigma :=&$\Sigma\ name\ E\ E\ |\ (E,E)\ |\ E.1\ |\ E.2\ |\ E.name$ \\
     id :=&$\textbf{Id}\ E\ |\ \textbf{ref}\ E\ |\ \textbf{id$_J$}\ E$ \\
      0 :=&$\textbf{0}\ |\ \textbf{ind$_0$}\ \textrm{E}\ \textrm{E}\ \textrm{E}$ \\
      1 :=&$\textbf{1}\ |\ \star\ |\ \textbf{ind$_1$}\ \textrm{E}\ \textrm{E}\ \textrm{E}$ \\
      2 :=&$\textbf{2}\ |\ 0_2\ |\ 1_2\ |\ \textbf{ind$_2$}\ \textrm{E}\ \textrm{E}\ \textrm{E}$ \\
      W :=&$\textbf{W}\ \textrm{ident}\ \textrm{E}\ \textrm{E}\ |\ \textbf{sup}\ \textrm{E}\ \textrm{E}\ |\ \textbf{ind$_W$}\ \textrm{E}\ \textrm{E}$ \\
\end{tabular}
\end{table}

\newpage
\subsection{Anders: Cubical HoTT \cite{cchm}}\footnote{\url{https://anders.groupoid.space}}
Semantically motivated by Coquand's Groupoid-Valued Presheaf models (where the Grothendieck construction models context extensions), Anders implements the CCHM cubical structure with Path, Glue, and higher inductive types (Coequ, Disc).
\begin{table}[ht!]
\centering
\begin{tabular}{rl}
 cosmos :=&$\textbf{U}_j \ |\ \textbf{V}_k$ \\
    var :=&$\textbf{var}\ name\ |\ \textbf{hole}$ \\
     pi :=&$\Pi\ name\ E\ E\ |\ \lambda\ name\ E\ E\ |\ E\ E$ \\
   sigma :=&$\Sigma\ name\ E\ E\ |\ (E,E)\ |\ E.1\ |\ E.2$ \\
      0 :=&$\textbf{0}\ |\ \textbf{ind$_0$}\ \textrm{E}\ \textrm{E}\ \textrm{E}$ \\
      1 :=&$\textbf{1}\ |\ \star\ |\ \textbf{ind$_1$}\ \textrm{E}\ \textrm{E}\ \textrm{E}$ \\
      2 :=&$\textbf{2}\ |\ 0_2\ |\ 1_2\ |\ \textbf{ind$_2$}\ \textrm{E}\ \textrm{E}\ \textrm{E}$ \\
      W :=&$\textbf{W}\ \textrm{ident}\ \textrm{E}\ \textrm{E}\ |\ \textbf{sup}\ \textrm{E}\ \textrm{E}\ |\ \textbf{ind$_W$}\ \textrm{E}\ \textrm{E}$ \\
     id :=&$\textbf{Id}\ E\ |\ \textbf{ref}\ E\ |\ \textbf{id$_J$}\ E$ \\
   path :=&$\textbf{Path}\ E\ |\ E^i\ |\ E\ @\ E$ \\
      I :=&$\textbf{I}\ |\ 0\ |\ 1\ |\ E\ \bigvee\ E\ |\ E\ \bigwedge\ E\ |\ \neg E$ \\
   part :=&$\textbf{Partial}\ E\ E\ |\ \textbf{[}\ (E=I) \rightarrow E, ...\ \textbf{]}$ \\
    sub :=&$\textbf{inc}\ E\ |\ \textbf{ouc}\ E\ |\ E\ \textbf{[}\ I\ \mapsto\ E\ \textbf{]}$ \\
    kan :=&$\textbf{transp}\ E\ E\ |\ \textbf{hcomp}\ E$ \\
   glue :=&$\textbf{Glue}\ E\ |\ \textbf{glue}\ E\ |\ \textbf{unglue}\ E\ E$ \\
  coequ :=&$\textbf{coequ}\ \mathrm{E}\ |\ \iota_2\ |\ \textbf{resp}\ |\ \textbf{ind}_{coequ}$ \\
   disc :=&$\textbf{disc}\ \mathrm{E}\ |\ \textbf{base}\ |\ \textbf{hub}\ |\ \textbf{spoke}\ |\ \textbf{ind}_{disc}$ \\
\end{tabular}
\end{table}

\subsection{Urs: The Family of Layered Sublanguages}\footnote{\url{https://urs.groupoid.space}}

\begin{remark}[The 3-Layer Architecture]
The sublanguages represent a progressive structure for formalizing mathematical and
physical concepts:
\begin{enumerate}
  \item \textbf{Felix}: Graded universes and tensors, group
    actions, and super-modality operators for geometric cohesion and
    differential structure---bosonic/fermionic grades, flat/sharp modalities.
  \item \textbf{Jack}: Homotopical foundations via
    spectra and groupoids, supporting differential K-theory with differential
    forms, connections, and refined cohomology structures.
  \item \textbf{Urs}: Configuration spaces and braids for Urs
    statistics, supporting Urs systems and linear types ($!A$) for
    resource-sensitive Urs programming.
\end{enumerate}
Each sublanguage culminates in a system tailored to formalize superpoints
($\mathbb{R}^{m|n}$), supersymmetry, and equivariant structures.
\end{remark}

\newpage
\subsubsection{Felix (Layer I) \cite{schreibershulman12cohesive}}\footnote{\url{https://felix.groupoid.space}}
Extends the homotopy core with graded universes, tensors, group actions, and cohesive modalities ($\flat$, $\sharp$, $\Im$) for synthetic supergeometry.
\begin{table}[ht!]
\centering
\begin{tabular}{rl}
   grade :=&$\textbf{Bose}\ |\ \textbf{Fermi}$ \\
  cosmos :=&$\textbf{U}_{j,g}$ \\
    smth :=&$\textbf{SmthSet}\ |\ \textbf{Plot}\ j\ E\ E\ |\ \textbf{SupSmthSet}$ \\
    cohs :=&$\flat\ E\ |\ \sharp\ E\ |\ \Im\ E$ \\
    bose :=&$\bigcirc\ E$ \\
  tensor :=&$E\ \otimes\ E$ \\
\end{tabular}
\end{table}

\subsubsection{Jack (Layer II) \cite{schreiber13differential}}\footnote{\url{https://jack.groupoid.space}}
Extends cohesion with stable homotopy (spectra, suspensions, wedges) and differential cohomology (forms, connections, and differential K-theory).
\begin{table}[ht!]
\centering
\begin{tabular}{rl}
    grpd :=&$\textbf{Grpd}\ j\ |\ \textbf{Comp}\ j\ E\ E\ E$ \\
  stable :=&$\textbf{Spectrum}\ |\ \textbf{Susp}\ E\ |\ E\ \wedge\ E\ |\ \textbf{Hom}_{Spec}\ E\ E$ \\
      ku :=&$KU_G^{\tau}\ E\ E\ E\ |\ \textbf{forms}\ j\ E\ |\ d\ j\ E\ |\ \textbf{DiffKU}_G\ E\ E\ E\ E$ \\
\end{tabular}
\end{table}

\subsubsection{Urs (Layer III) \cite{sati2022anyonic, sati2022topological, sati2023anyonicdefect}}\footnote{\url{https://urs.groupoid.space}}
The top layer, adding configuration spaces, braids, and linear resource types ($!A$) for Urs programming.
\begin{table}[ht!]
\centering
\begin{tabular}{rl}
   qubit :=&$\textbf{Qubit}\ E\ E\ |\ \textbf{app}\ E\ E\ E\ |\ \textbf{fuse}\ E\ E\ E$ \\
  linear :=&$!\ E$ \\
    topo :=&$\textbf{Config}_j\ E\ |\ \textbf{Braid}_j\ E$ \\
\end{tabular}
\end{table}

\begin{remark}[STT $\to$ Urs: The Simplicial Connection]
Simplicial type theory (Ulrik/STT) provides a direct connection to Urs.
The directed interval $\mathbb{I}^{\to}$, extension types, and modal
$\Pi$ in STT supply the $\infty$-categorical framework that Urs's Layer I
(Homotopy) builds upon. By isolating the category of categories $\mathbf{Cat}$ as a reflective subuniverse in simplicial type theory, STT provides the required directed homotopy structures. Concretely, STT's directed interval becomes the path/homotopy backbone, while its modal/shape operations inform Urs's cohesive modalities ($\flat$, $\sharp$, $\Im$) in Layer II\@. This is represented by the $\mathrm{STT}$ arrow from Ulrik to Urs in the lattice diagram, which is \emph{not} simply a factoring through Anders.
\end{remark}

\newpage
\section{Spine III: Infinity Category Theory}
This spine encompasses operational algebra and directed category theories, dealing with presentations of algebraic shapes and directed categorical structures.

\subsection{Dan: Operational Algebra \cite{birkedal2011first}}\footnote{\url{https://dan.groupoid.space}}
Dan operates on a completely different level: it manipulates algebraic
presentations (groups, simplicial sets, chain complexes) operationally.
Its term language has no dependent types, no universes, and no $\Pi$/$\Sigma$.
Instead it works with group words, permutations, matrices, face/degeneracy
maps, and boundary operators.

\begin{table}[ht!]
\centering
\begin{tabular}{rl}
    type :=&$\textbf{Simplex}\ |\ \textbf{Group}\ |\ \textbf{Simplicial}\ |\ \textbf{Chain}\ |\ \textbf{Category}$ \\
         &$\textbf{Monoid}\ |\ \textbf{Ring}\ |\ \textbf{Field}\ |\ \textbf{PermGroup}\ |\ \textbf{MatGroup}$ \\
         &$\textbf{FpGroup}\ |\ \textbf{PcGroup}\ |\ \textbf{Module}\ |\ \textbf{Set}\ |\ \textbf{GSet}$ \\
    term :=&$name\ |\ E\ \circ\ E\ |\ E^{-1}\ |\ E^j\ |\ e\ |\ \textbf{matrix}\ \textrm{args}$ \\
         &$E\ +\ E\ |\ E\ \cdot\ E\ |\ \textbf{perm}\ \textrm{args}\ |\ \textbf{cycle}\ \textrm{args}$ \\
         &$\textbf{PermGroup}\ \textrm{args}\ |\ \textbf{MatGroup}\ \textrm{args}\ j\ k$ \\
         &$\textbf{FpGroup}\ \textrm{names}\ \textrm{args}\ |\ \textbf{PcGroup}\ \textrm{names}\ \textrm{args}$ \\
         &$d_j\ E\ |\ s_j\ E\ |\ \partial\ E\ |\ E\ \cdot\ E\ |\ \textbf{Hom}\ E\ E\ |\ E\ \equiv\ E$ \\
    hypo :=&$\textbf{decl}\ \textrm{names}\ type\ |\ \textbf{eq}\ name\ E\ E\ |\ \textbf{map}\ name\ E\ E$ \\
         &$\textbf{pres}\ name\ E\ |\ \textbf{rel}\ E\ |\ \textbf{act}\ name\ E\ E$ \\
         &$\textbf{orbit}\ E\ E\ \textrm{args}\ |\ \textbf{stabilizer}\ E\ E\ E$ \\
\end{tabular}
\end{table}

\subsection{Mike: Directed 1-Category Type Theory \cite{Shulman15}}\footnote{\url{https://mike.groupoid.space}}
A \emph{linear} type theory with a bidirectional type checker enforcing
quadraticality (each category variable occurs exactly once covariantly and
once contravariantly).
\begin{table}[ht!]
\centering
\begin{tabular}{rl}
     var :=&$\textbf{var}\ name\ |\ name\ \textrm{args}$ \\
      op :=&$\textbf{op}\ E$ \\
    prod :=&$E\ \times\ E$ \\
     hom :=&$\textbf{Hom}_E\ E\ E$ \\
  tensor :=&$E\ \otimes\ E$ \\
      pi :=&$E\ \multimap\ E\ |\ \lambda\ name\ E\ |\ E\ E$ \\
   coend :=&$\int^{E}\ name\ E\ |\ \mathbf{coend}\ name\ name\ name\ E$ \\
     end :=&$\int_{E}\ name\ E\ |\ \mathbf{end}\ name\ E$ \\
      id :=&$\textbf{id}\ E$ \\
     idJ :=&$\textbf{id}_J\ E\ |\ \textbf{id}_{J,cov}\ E\ |\ \textbf{id}_{J,contra}\ E$ \\
\end{tabular}
\end{table}

\newpage
\subsection{Ulrik: Simplicial \texorpdfstring{$\infty$}{Infinity}-Categories \cite{riehl2017type, gratzer2025yoneda}}\footnote{\url{https://ulrik.groupoid.space}}
Extends Mike's foundation with directed interval $\mathbb{I}$, shape
inclusions ($\varphi \subseteq \psi$), extension types, modal $\Pi$, and
twisted arrow categories ($A^{\mathrm{tw}}$).

To construct a type theory for synthetic category theory, one might hope to interpret type theory in the category of categories ($\infty$-categories or otherwise) to ensure that types realize categories. However, the category of small categories is too poorly behaved to form a model of Martin-Löf Type Theory (MLTT). Instead, Riehl and Shulman (2017) \cite{riehl2017type} use the category of bisimplicial sets (simplicial spaces) to model synthetic $(\infty, 1)$-categories via Rezk spaces, enabling further formalizations such as the Yoneda lemma and presheaf categories \cite{gratzer2025yoneda}. Simplicial Type Theory (STT) axiomatizes this framework using a directed interval and shapes, allowing one to isolate $(\infty, 1)$-categories as a reflective subuniverse of types within the theory.

\begin{table}[ht!]
\centering
\begin{tabular}{rl}
 cosmos :=&$\mathbf{U}$ \\
    var :=&$\textbf{var}\ name$ \\
     pi :=&$\Pi\ name\ E\ E\ |\ \lambda\ name\ E\ E\ |\ E\ E$ \\
  sigma :=&$\Sigma\ name\ E\ E\ |\ (E,E)\ |\ E.1\ |\ E.2$ \\
     id :=&$\textbf{Id}\ E\ E\ E\ |\ \textbf{ref}\ E$ \\
   idir :=&$\mathbb{I}\ |\ 0\ |\ 1\ |\ E\ \leq\ E\ |\ \Phi\ \subseteq\ \Psi$ \\
   part :=&$\textbf{Partial}\ E\ E\ |\ \textbf{[}\ (E=I) \rightarrow E, ...\ \textbf{]}$ \\
    ext :=&$\textbf{Ext}\ E\ E\ E$ \\
 modalpi :=&$\Pi\ \textrm{shape}\ E\ E\ |\ \lambda\ \textrm{shape}\ E\ |\ E\ E$ \\
     tw :=&$\textbf{Tw}\ E\ |\ \textbf{tw}_{p0}\ E\ |\ \textbf{tw}_{p1}\ E$ \\
    ops :=&$E\ \vee\ E\ |\ E\ \wedge\ E\ |\ \neg E$ \\
     idJ :=&$\textbf{id$_J$}\ E\ |\ \textbf{id$_{J,cov}$}\ E\ |\ \textbf{id$_{J,contra}$}\ E$ \\
\end{tabular}
\end{table}

\newpage
\section{Cohomology of Programming Languages}
Each language's syntax is analyzed \emph{individually} as a CW-complex.

For a formal language $\mathcal{O}_x$, we associate a chain complex of abelian groups (free $\mathbb{Z}$-modules) $C_*(\mathcal{O}_x)$ constructed from the syntax and reduction/conversion semantics. This construction models the syntactic space as a cell complex (CW-complex) where paths and higher homotopies encode the computational rules.

\begin{definition}[Syntactic CW-Complex]
For a language $\mathcal{O}_x$ with term syntax $T$ generated by a signature $\Sigma$, the syntactic cell complex $C_*(\mathcal{O}_x)$ is defined in low dimensions by:
\begin{itemize}
    \item \textbf{0-cells
    ($C_0(\mathcal{O}_x)$) }: The free $\mathbb{Z}$-module generated by the set of basic syntactic
      constructors $c \in \Sigma$ (e.g., \texttt{Var}, \texttt{Pi}, \texttt{Lam}, \texttt{App} for Henk).
      Any term $t \in T$ has an associated constructor-occurrence representation:
    \[
    [t] = \sum_{c \in \Sigma} \mathrm{count}(c, t) \cdot c \;\; \in C_0(\mathcal{O}_x)
    \]
    where $\mathrm{count}(c, t)$ is the number of times the constructor $c$ appears in the syntax tree of $t$.
    \item \textbf{1-cells
    ($C_1(\mathcal{O}_x)$)}: The free $\mathbb{Z}$-module generated by the rewriting and reduction rules
    $r \in \mathcal{R}$ of the language (e.g., $\beta$-reductions, evaluation steps). Each $1$-cell $r : t \to t'$
    represents a directed rewrite path from a redex term $t$ to a contractum term $t'$.
    \item \textbf{2-cells
    ($C_2(\mathcal{O}_x)$)}: The free $\mathbb{Z}$-module generated by equivalence relations, critical pairs,
    confluence diagrams, and conversion equations (e.g., $\eta$-conversions, commuting conversions). A $2$-cell
    $e$ represents a coherence condition or a closed loop of rewrite steps.
    \item \textbf{Higher cells
    ($C_k(\mathcal{O}_x)$ for $k \geq 3$)}: Coherence relations between equations, such as Mac Lane pentagon
    relations, path-type univalence coherences (in Anders), or quadraticality constraints (in Mike).
\end{itemize}
\end{definition}

\begin{definition}[Boundary Operators]
The boundary operators $\partial_k : C_k(\mathcal{O}_x) \to C_{k-1}(\mathcal{O}_x)$ are defined as follows:
\begin{enumerate}
    \item The 1-boundary operator $\partial_1 : C_1(\mathcal{O}_x) \to C_0(\mathcal{O}_x)$ is the linear map determined by the net change in constructor counts caused by a rewrite rule $r : t \to t'$:
    \[ \partial_1(r) = [t'] - [t] \]
    \item The 2-boundary operator $\partial_2 : C_2(\mathcal{O}_x) \to C_1(\mathcal{O}_x)$ is determined by the boundary of confluence diagrams or conversion loops. If a 2-cell $e$ represents a confluence diagram from a term $s$ reducing to $s'$ via two different paths of rewrites $p_1 = (r_{1,1}, \dots, r_{1,m})$ and $p_2 = (r_{2,1}, \dots, r_{2,n})$, the boundary is:
    \[ \partial_2(e) = \sum_{j=1}^m r_{1,j} - \sum_{k=1}^n r_{2,k} \]
    Since both paths $p_1$ and $p_2$ start at $s$ and end at $s'$, the net constructor changes must match:
    \[ \partial_1(p_1) = [s'] - [s] = \partial_1(p_2) \]
    Thus:
    \[ \partial_1(\partial_2(e)) = \partial_1(p_1) - \partial_1(p_2) = 0 \]
    This proves that $\partial_1 \circ \partial_2 = 0$, establishing a valid chain complex.
\end{enumerate}
\end{definition}

\begin{definition}[Syntactic Cohomology]
Let $C_*(\mathcal{O}_x)$ be the syntactic chain complex of $\mathcal{O}_x$ and let $G$ be an abelian group of coefficients. The cochain complex $C^*(\mathcal{O}_x; G) = \mathrm{Hom}(C_*(\mathcal{O}_x), G)$ has coboundary operators $\delta^k : C^k(\mathcal{O}_x; G) \to C^{k+1}(\mathcal{O}_x; G)$ defined by $\delta^k(\phi) = \phi \circ \partial_{k+1}$.
The syntactic cohomology groups are:
\[ H^k(\mathcal{O}_x; G) = \ker(\delta^k) / \mathrm{im}(\delta^{k-1}) \]
\end{definition}

\subsection{Physical Interpretations of Cohomology Groups}
The cohomology groups $H^k(\mathcal{O}_x; G)$ measure fundamental properties of the syntax and semantics of $\mathcal{O}_x$:
\begin{itemize}
    \item $H^0(\mathcal{O}_x; G)$: Measures the \emph{syntactic dimension} or independent constructor classes. A class in $H^0$ is a function $f : C_0 \to G$ such that for any rewrite $r : t \to t'$, $f([t]) = f([t'])$. Thus, $H^0$ captures invariants preserved under computation (such as type preservation or conserved weights of terms).
    \item $H^1(\mathcal{O}_x; G)$: Measures \emph{normalization obstructions}. An element in $H^1$ is represented by a 1-cocycle (an assignment of weights to rewrites that sums to zero on closed loops) modulo those coming from constructor valuations. Non-trivial $H^1$ classes indicate non-confluent rewrites, cycles of reductions (non-termination), or irreversible compilation steps.
    \item $H^2(\mathcal{O}_x; G)$: Measures \emph{coherence obstructions}. It captures the obstructions to filling confluence diagrams with higher-dimensional conversion cells. A non-trivial class in $H^2$ represents a critical pair (or confluence loop) that cannot be closed by conversions or equivalence cells, pointing to deep semantic mismatches or failures of univalence/coherence.
\end{itemize}

\newpage
\subsection{Examples across the Spines}

\begin{example}[Spines I/II: Strict \& Homotopical]
\end{example}
\begin{itemize}
    \item \textbf{Henk
    ($\mathcal{O}_\Pi$)}: The core type theory.
    \[ C_0 = \{ \mathtt{Var}, \mathtt{Universe}, \mathtt{Pi}, \mathtt{Lam}, \mathtt{App} \} \]
    The 1-cells are generated by $\beta$-reduction:
    \[ r_\beta : \mathtt{App}(\mathtt{Lam}(x, A, t), u) \to t[u/x] \]
    $\partial_1(r_\beta) = [t[u/x]] - ( \mathtt{App} + \mathtt{Lam} + [t] + [u] )$.
    The 2-cells consist of $\eta$-conversions:
    \[ e_\eta : \mathtt{Lam}(x, A, \mathtt{App}(t, x)) \equiv_\eta t \]
    Since Henk is strongly normalizing and confluent (Church-Rosser), the higher cohomology $H^k(\mathcal{O}_\Pi; \mathbb{Z})$ for $k \geq 1$ is trivial, reflecting a contractible syntactic space.
    \item \textbf{Anders
    ($\mathcal{O}_{\mathrm{Anders}}$)}: Adds univalence and Glue types. The univalence axiom represents a path-space equivalence between isomorphic types. This introduces higher-dimensional coherences (Glue cells) acting as 3-cells, preventing non-trivial 2-dimensional obstructions.
    \item \textbf{Urs
    ($\mathcal{O}_{\mathrm{Urs}}$)}: Adds modalities ($\flat, \sharp$) and differential structures (forms, connections). The syntactic CW-complex of $\mathcal{O}_{\mathrm{Urs}}$ is enriched with topological and differential cells, yielding non-trivial $H^k$ corresponding to de Rham cohomology and Chern-Weil invariants of smooth/differential type structures.
\end{itemize}

\begin{example}[Spine III: Category Theory]
\end{example}
\begin{itemize}
    \item \textbf{Dan
    ($\mathcal{O}_{\mathrm{Dan}}$)}: Operational algebra. Here, term syntax encodes algebraic operations (e.g., face maps $\partial_i$, degeneracies $s_i$). The cell complex $C_*(\mathcal{O}_{\mathrm{Dan}})$ is isomorphic to the standard simplicial chain complex of the underlying category/algebraic presentation.
    \item \textbf{Mike
    ($\mathcal{O}_{\mathrm{Mike}}$)}: Directed category theory. 1-cells are non-reversible functors/morphisms, meaning paths are directed. The boundary operator $\partial_1$ tracks directed change. Quadraticality constraints act as 2-cells that model the non-commutative relation profiles of monoidal/directed types.
    \item \textbf{Ulrik
    ($\mathcal{O}_{\mathrm{Ulrik}}$)}: Simplicial $\infty$-categories. Generalizes Mike's 1-categories to directed simplicial spaces. The higher cells ($C_k$ for $k \geq 2$) represent the Segal completeness conditions and Rezk completeness conditions, ensuring that the directed paths compose coherently up to higher homotopies.
\end{itemize}

\newpage
\section{Spectral Sequence of Language Categories}
Unlike a simple linear filtration, the language category $\mathcal{L}$ is
organized as a \emph{lattice} (a partially ordered set $\mathcal{I}$) with multiple join-paths. 
To compute the cohomology of the total language $\mathcal{L}$ from its sublanguages, we can employ the Bousfield-Kan spectral sequence for the homotopy colimit over the poset category $\mathcal{I}$:
\[ E_2^{p,q} = H^p(\mathcal{I}; \mathcal{H}^q) \implies H^{p+q}(\mathcal{L}) \]
where $\mathcal{H}^q : \mathcal{I} \to \mathbf{Ab}$ is the functor mapping each language $\mathcal{O}$ to its cohomology $H^q(\mathcal{O})$, and $H^p(\mathcal{I}; -)$ denotes the cohomology of the category $\mathcal{I}$.

Alternatively, one can impose an integer rank filtration on the syntactic complex $C^*(\mathcal{L})$ by defining the depth $p$ of a language in the lattice. Let $F_p C^*(\mathcal{L})$ be the subcomplex generated by all languages of depth $\le p$. The associated spectral sequence is:
\[ E_1^{p,q} = H^{p+q}(F_p / F_{p-1}) \implies H^{p+q}(\mathcal{L}) \]
where the relative complex $F_p / F_{p-1}$ isolates the ``new'' syntactic primitives introduced exactly at depth $p$.

\begin{table}[ht]
\centering
\resizebox{\textwidth}{!}{%
\begin{tabular}{c|l|l|l|l}
\hline
\textbf{$p$} & \textbf{Language} & \textbf{Branch} & \textbf{Primitives added}
& \textbf{$|C_0|$} \\
\hline
0 & Henk & CTT & $\Pi$, $U^n$ & 5 \\
1 & Frank & CTT & + Ind & 8 \\
2a & Christine & CTT & + $\Sigma$, Id & 14 \\
2b & Per & CTT & + $\Sigma$, Id, W, 0, 1, 2 & 26 \\
3 & Anders & CTT & + Path, I, Glue, HComp, HIT & 40+ \\
4a & Felix & CTT/Modal & + $\flat$, $\sharp$, $\Im$, $\otimes$, Plot & 10 \\
4b & Jack & CTT/Modal & + Spectra, Susp, Grpd, Forms, DiffKU & 12 \\
4c & Urs & CTT/Modal & + Qubit, Linear, Config, Braid & 7 \\
\hline
--- & Dan & Algebra & Groups, Simplices, Chains, Boundaries & 20 \\
--- & Mike & Directed & Hom, $\otimes$, $\int$, $\prod$, $\multimap$ & 13 \\
--- & Ulrik & Directed & + $\mathbb{I}^{\to}$, Ext, Tw, $\leq$ & 25 \\
\hline
\end{tabular}%
}
\caption{The language lattice: filtration indices and constructor counts.
Note the split at $p=2$ (Christine/Per), the four sublanguages of Urs at $p=4$ (Homotopy, SupSmth, DiffK, Urs), and the separate branch indices
for Dan, Mike, and Ulrik.}
\end{table}

\newpage

\begin{definition}[Morphisms in $\mathcal{L}$]
Morphisms $f : \mathcal{O}_x \to \mathcal{O}_y$ are classified as:
\begin{enumerate}
    \item $\mathcal{F}_{\mathrm{enrich}}$: Enrichment (type inference, elaboration).
    \item $\mathcal{F}_{\mathrm{simp}}$: Simplification (type erasure, extraction).
    \item $\mathcal{F}_{\mathrm{trans}}$: Translation (compilation across branches).
    \item $\mathcal{F}_{\mathrm{eval}}$: Evaluation (normalization to values).
    \item $\mathcal{F}_{\mathrm{norm}}$: Normalization ($\beta\eta$-reduction).
    \item $\mathcal{F}_{\mathrm{certify}}$: Certification (proof checking).
\end{enumerate}
Cross-branch morphisms (e.g.\ \textbf{Lift}: Dan $\to$ Mike, \textbf{Core}:
Ulrik $\to$ Anders) are structure-changing and do not preserve the term type.
\end{definition}

\subsection{Functorial Conversions and Cross-Branch Morphisms}
The language category $\mathcal{L}$ is not merely a set of objects;
the relationships between the branches are mediated by structured functorial conversions.
The primary conversions, their statuses, and associated information losses are detailed in Table~\ref{tab:functors}.

\begin{table}[ht]
\centering
\resizebox{\textwidth}{!}{%
\begin{tabular}{|l|c|c|l|p{8cm}|}
\hline
\textbf{Functor / Conversion} & \textbf{Source} & \textbf{Target} & \textbf{Status} & \textbf{Semantics \& Information Loss} \\
\hline
\textbf{Lift} & Dan & Mike & Yes & Semantical lifting of presentation generators and equations to 1-categorical structures. \\
\hline
\textbf{Forget} & Mike & Dan & Partial & Truncates categories back to operation-based presentation generators. \\
\hline
\textbf{Free} & Mike & Ulrik & Yes & Fully faithful embedding of directed 1-categories into directed $\infty$-categories. \\
\hline
\textbf{Trunc} & Ulrik & Mike & No & Cannot be done functorially without losing higher-dimensional coherence. \\
\hline
\textbf{Core} & Ulrik & Anders & Yes & Groupoid Core: discards directed path structure, retaining only the sub-groupoid of isomorphisms. \\
\hline
\textbf{Disc} & Anders & Ulrik & Yes & Discrete embedding: interprets $\infty$-groupoids as discrete $\infty$-categories via disc modality. \\
\hline
\textbf{Forget} & Anders & Mike & Partial & 1-category approximation: projects higher paths onto 1-category hom-sets. \\
\hline
\textbf{Groupoid} & Mike & Anders & Yes & Forgetful functor from directed 1-categories to spaces of isomorphism paths. \\
\hline
\textbf{Lift} (Composite) & Dan & Ulrik & Yes & Composition of \textbf{Lift} and \textbf{Free}: maps operations to Rezk $\infty$-categories. \\
\hline
\textbf{Forget} & Ulrik & Dan & No & Directed simplicial spaces are too rich to map back to discrete operations directly. \\
\hline
\textbf{Forget} & Anders & Dan & No & Undirected spaces cannot be functorially mapped to discrete presentation equations. \\
\hline
\textbf{Modal} & Anders & Felix & Yes & Extends cubical HoTT with cohesive modalities ($\flat$, $\sharp$, $\Im$) and super-grades. \\
\hline
\textbf{Diff} & Felix & Jack & Yes & Stable extension to spectra, suspensions, forms, and differential K-theory. \\
\hline
\textbf{Urs} & Jack & Urs & Yes & Extends differential K-theory with configuration spaces, braids, and linear types. \\
\hline
\end{tabular}%
}
\caption{Primary conversions and functors between the language layers.}
\label{tab:functors}
\end{table}

\newpage

\subsubsection{Deep Dive: The Lift Functor
(\texorpdfstring{$\mathbf{Dan} \to \mathbf{Ulrik}$}{Dan -> Ulrik})}

The \textbf{Lift} functor acts as a semantical compilation bridge that transports combinatorial presentations from
the operational layer $\mathbf{Dan}$ to the directed synthetic type theories of $\mathbf{Mike}$ and $\mathbf{Ulrik}$.
Rather than a simple syntax translator, it elaborates discrete presentations into rich categorical models in four distinct phases:

\begin{enumerate}
    \item \textbf{Signature Ingestion}: Reads the generator set, face maps ($\partial_i$), degeneracy maps, and algebraic presentation relations directly from the $\mathbf{Dan}$ syntax.
    \item \textbf{Semantical Lifting}: Maps the ingested signature into category-theoretic structures, translating generator equations into functors, categories, and sheaves.
    \item \textbf{Elaboration to STT}: Translates composition equations to identity paths ($\mathtt{EId}$) in $\mathbf{Ulrik}$. The directed interval $\mathbb{I}^{\to}$ is employed to model the domains and boundary conditions.
    \item \textbf{Metatheory Absorption}: Automatically equips the lifted object with the abstract category-theoretic library of $\mathbf{Ulrik}$ (including Yoneda lemmas, limits, adjunctions, and Kan extensions).
\end{enumerate}

\subsubsection{Deep Dive: The Groupoid Core (\texorpdfstring{$\mathbf{Ulrik} \to \mathbf{Anders}$}{Ulrik -> Anders})}
The translation from $\mathbf{Ulrik}$ to $\mathbf{Anders}$ is a structural forgetful conversion that extracts the underlying homotopy type (the $\infty$-groupoid of isomorphisms) from a directed $\infty$-category. Since the simplicial directed interval $\mathbb{I}^{\to}$ and directed hom-types ($x \to y$) of $\mathbf{Ulrik}$ do not exist in the cubical undirected world of $\mathbf{Anders}$, the functor operates as follows:
\begin{itemize}
    \item Discards directed arrows, keeping only the sub-groupoid $A^{\mathrm{core}}$ containing morphisms with valid inverses.
    \item Discards the Segal and Rezk completeness conditions, transforming directed category composition into symmetric path composition.
    \item Replaces the directed interval $\mathbb{I}^{\to}$ with the cubical symmetric interval $I$ (with connection and reversal operations).
\end{itemize}

\newpage
\subsubsection{Case Study: M\"obius Strip Representation}
To illustrate the conversions, Table~\ref{tab:mobius} compares the representation of a M\"obius Strip across the three primary layers ($\mathbf{Dan}$, $\mathbf{Ulrik}$, and $\mathbf{Anders}$).

\begin{table}[ht]
\centering
\caption{Comparison of the M\"obius Strip representation across the layers.}
\label{tab:mobius}
\noindent\makebox[\textwidth][c]{%
\begin{tabular}{llll}
\hline
\textbf{Aspect} & \textbf{Dan (Operational)} & \textbf{Ulrik (Simplicial STT)} & \textbf{Anders (Cubical HoTT)} \\
\hline
\textbf{Framework} & Algebra presentations & STT with directed interval $\mathbb{I}^{\to}$ & Cubical HoTT with modalities \\
\textbf{Math Type} & Simplicial complex with boundaries & Rezk $\infty$-category from simplices & Bundle $\sum_{x: S^1} \mathrm{Moebius}(x)$ \\
\textbf{Twist Encoding} & Equations on face maps ($\partial_i$) & Directed paths satisfying Segal & Path transport ($\mathtt{ua}$) via $\mathtt{not}$ \\
\textbf{Homotopy Type} & Contractible local complex ($\ast$) & $\infty$-category with core $\simeq S^1$ & Homotopy equivalence to $S^1$ \\
\textbf{Abstraction} & Combinatorial building blocks & Synthetic directed $\infty$-category & Synthetic univalent bundle \\
\textbf{Functor Action} & Ingested by \textbf{Lift} & Lifted to directed $\infty$-category & Extracted via \textbf{Core} \\
\textbf{Information} & Combinatorial generators & Full directed structure & Undirected homotopy skeleton \\
\hline
\end{tabular}%
}
\end{table}

\section{Conclusion}

A recurring ambition in the development of higher-dimensional type theories is the pursuit of a singular, comprehensive foundational language---a formal system capable of expressing discrete combinatorial presentations, directed synthetic category theory, and univalent cohesive geometry simultaneously. However, as demonstrated by the categorical obstructions documented in this article, we conclude that \emph{no natural unified foundation is viable}. Any attempt to internalize this entire hierarchy into a single categorical semantics would inevitably require rigidification into strict bicategories, tricategories, or infinitely coherent higher-categorical structures. Such a global strictification is highly non-natural and introduces prohibitive complexity. Instead, the mathematical reality of higher-dimensional formalisms naturally forms a \emph{functorial mesh of transports}.

The obstructions to a unified foundation are deeply rooted in the underlying geometry of the respective topoi. For example, consider the embedding of Simplicial Type Theory (\textbf{Ulrik}) into Cubical HoTT (\textbf{Anders}). The directed interval $\mathbb{I}^{\to}$ and its associated Segal and Rezk completeness conditions depend fundamentally on an asymmetrical, directed path space. When attempting to interpret this within the symmetric, invertible cubical paths of \textbf{Anders}, one is forced to invoke the \textbf{Core} functor. This functor aggressively truncates the directed $\infty$-category, extracting only its maximal $\infty$-groupoid of isomorphisms. The directed geometric structure is irreversibly collapsed by this essential geometric morphism.

Similarly, the transition from discrete algebraic presentations (\textbf{Dan}) to synthetic homotopy theories cannot be accomplished via internal deformation. The combinatorial algebra of face maps, degeneracy maps, and strict boundary operators resides in a category devoid of continuous path spaces or univalent universes. To interpret these discrete generators within the synthetic structures of \textbf{Mike} or \textbf{Ulrik}, one must employ the \textbf{Lift} functor to fully realize them as categorical objects. There exists no continuous adjunction from this rigid, discrete combinatorial layer into the cohesive $\infty$-topoi of \textbf{Felix} or the linear configuration spaces of \textbf{Urs}.

Ultimately, the architecture of modern higher-dimensional mathematics is irreducibly pluralistic. The ``Groupoid Infinity'' stack acknowledges this by treating specialized theories not as fragments of a missing monolith, but as distinct objects in a category of formal languages. By relying on a functorial mesh of geometric morphisms, discrete embeddings, and modal coreflections, we achieve a robust and mathematically rigorous ecosystem whose compositional power far exceeds that of any artificially strictified unified theory.

\begin{thebibliography}{CCHM17}

\bibitem[AKS15]{Shulman15}
Benedikt Ahrens, Krzysztof Kapulkin, and Michael Shulman.
\newblock Univalent categories and the rezk completion.
\newblock In Maria del~Mar Gonz{\'a}lez, Paul~C. Yang, Nicola Gambino, and
  Joachim Kock, editors, {\em Extended Abstracts Fall 2013}, pages 75--76,
  Cham, 2015. Springer International Publishing.

\bibitem[Bar92]{Henk93}
H.~P. Barendregt.
\newblock Lambda calculi with types.
\newblock In S.~Abramsky, Dov~M. Gabbay, and S.~E. Maibaum, editors, {\em
  Handbook of Logic in Computer Science (Vol. 2)}, pages 117--309, New York,
  NY, USA, 1992. Oxford University Press, Inc.

\bibitem[BMSS11]{birkedal2011first}
Lars Birkedal, Rasmus~Ejlers Mogelberg, Jan Schwinghammer, and Kristian
  Stovring.
\newblock First steps in synthetic guarded domain theory: step-indexing in the
  topos of trees.
\newblock In {\em 2011 IEEE 26th Annual Symposium on Logic in Computer
  Science}, pages 55--64. IEEE, 2011.

\bibitem[CCHM17]{cchm}
Cyril Cohen, Thierry Coquand, Simon Huber, and Anders M{\"{o}}rtberg.
\newblock Cubical type theory: {A} constructive interpretation of the
  univalence axiom.
\newblock {\em {FLAP}}, 4(10):3127--3170, 2017.

\bibitem[CH88]{Coq88}
Thierry Coquand and Gerard Huet.
\newblock The calculus of constructions.
\newblock In {\em Information and Computation}, pages 95--120, Duluth, MN, USA,
  1988. Academic Press, Inc.

\bibitem[Sch50]{schwartz50distributions}
Laurent Schwartz.
\newblock {\em Th{\'e}orie des distributions}.
\newblock Hermann, Paris, 1950.

\bibitem[MLS84]{Lof84}
P.~Martin-L{\"o}f and G.~Sambin.
\newblock {\em Intuitionistic type theory}.
\newblock Studies in proof theory. Bibliopolis, 1984.

\bibitem[Pau93]{paulinmohring}
Christine Paulin{-}Mohring.
\newblock Inductive definitions in the system {C}oq --- rules and properties.
\newblock In {\em Typed Lambda Calculi and Applications, International
  Conference on Typed Lambda Calculi and Applications, {TLCA} '93, Utrecht, The
  Netherlands, March 16-18, 1993, Proceedings}, pages 328--345, 1993.

\bibitem[PP89]{Pfenning89}
Frank Pfenning and Christine Paulin{-}Mohring.
\newblock Inductively defined types in the calculus of constructions.
\newblock In {\em Mathematical Foundations of Programming Semantics, 5th
  International Conference, Tulane University, New Orleans, Louisiana, USA,
  March 29 - April 1, 1989, Proceedings}, pages 209--228, 1989.

\bibitem[GWB25]{gratzer2025yoneda}
Daniel Gratzer, Jonathan Weinberger, and Ulrik Buchholtz.
\newblock The Yoneda embedding in simplicial type theory.
\newblock Ukrainian Translation: \url{https://groupoid.space/friends/ulrik/yoneda-simplicialtt-ua.pdf}, 2025.

\bibitem[RS17]{riehl2017type}
Emily Riehl and Michael Shulman.
\newblock A type theory for synthetic $\infty$-categories.
\newblock {\em arXiv preprint arXiv:1705.07442}, 2017.

\bibitem[Sch13]{schreiber13differential}
Urs Schreiber.
\newblock {Differential cohomology in a cohesive infinity-topos}.
\newblock 2013.

\bibitem[SS12]{schreibershulman12cohesive}
Urs Schreiber and Michael Shulman.
\newblock Quantum gauge field theory in cohesive homotopy type theory.
\newblock In {\em Workshop on Quantum Physics and Logic}, 2012.

\bibitem[SS22a]{sati2022anyonic}
Hisham Sati and Urs Schreiber.
\newblock Anyonic topological order in twisted equivariant differential (TED) K-theory.
\newblock {\em arXiv preprint arXiv:2206.13563}, 2022.

\bibitem[SS22b]{sati2022topological}
Hisham Sati and Urs Schreiber.
\newblock Topological quantum programming in TED-K.
\newblock {\em arXiv preprint arXiv:2209.08331}, 2022.

\bibitem[SS23]{sati2023anyonicdefect}
Hisham Sati and Urs Schreiber.
\newblock Anyonic defect branes and conformal blocks in twisted equivariant differential (TED) K-theory.
\newblock {\em Reviews in Mathematical Physics}, 35(08):2350024, 2023.

\end{thebibliography}
\end{document}